\chapter*{Kết luận}
\addcontentsline{toc}{chapter}{Kết luận}

\textbf{Kết quả đạt được.} Báo cáo đã tìm hiểu bài báo của Kassem~\cite{kassem2022} về quản lý rủi ro trong dự án xây dựng dầu khí bằng PLS-SEM và thu được các kết quả sau:
\begin{enumerate}
  \item \textbf{Hệ thống hóa kiến thức.} Báo cáo trình bày khái niệm rủi ro và quy trình quản lý rủi ro theo PMBOK và ISO~31000, đồng thời sắp xếp lại 22 kỹ thuật quản lý rủi ro theo quy trình mà chúng phục vụ. Ba phương pháp RII, ma trận xác suất -- tác động và PLS-SEM được trình bày cùng công thức, thuật toán, ngưỡng đánh giá, ưu điểm, hạn chế và xu hướng.
  \item \textbf{Tái hiện thực nghiệm theo hai hướng.} Kiểm tra lại số liệu đã công bố cho thấy phần tính toán cơ bản của bài báo đúng: RII, AVE, CR, $Q^2$, và các chỉ số $R^2$, $f^2$ nhất quán với ma trận tương quan. Tuy nhiên, việc kiểm tra cũng phát hiện tiêu chí Fornell--Larcker không đạt ở nhóm RM, AVE của biến kết quả không khớp với loading, và không tính lại được AVE dùng trong GoF. Trên dữ liệu mô phỏng, $Q^2$ và sai số chuẩn của các đường dẫn bậc 2 trùng gần như tuyệt đối với bài báo. Ngược lại, hệ số và giá trị $t$ của 13 giả thuyết không tái hiện được, và HTMT vượt ngưỡng ở ba cặp nhóm rủi ro.
  \item \textbf{Liên hệ thực tế.} Hai trường hợp Nhiệt điện Thái Bình 2 và Healthcare.gov cho thấy khung 13 nhóm rủi ro của bài báo dùng lại được. Tuy vậy, ở bối cảnh ổn định thì rủi ro nội bộ (chủ đầu tư, nhà thầu, đấu thầu, quản lý dự án) mới là nguyên nhân chính, ngược với kết quả ở Yemen.
  \item \textbf{Áp dụng công cụ.} Báo cáo đã dùng MS Project, Excel và mô phỏng Monte Carlo trên một dự án hệ thống thông tin mẫu. Kết quả cho thấy phân tích định lượng bổ sung cho ma trận PIM: nó chỉ ra rủi ro nào thực sự kéo dài dự án và dự án cần dự phòng bao nhiêu thời gian. Các công cụ này cũng được dùng để quản lý chính quá trình làm báo cáo.
\end{enumerate}

\textbf{Đánh giá chung về bài báo.} Kết luận định tính chính của bài báo đứng vững: ở Yemen, rủi ro bên ngoài ảnh hưởng tới thành công của dự án mạnh hơn rủi ro nội bộ. Các kết luận chi tiết về từng nhóm rủi ro thì cần được xem xét thận trọng, do thiếu minh bạch về phương pháp và có nhiều mâu thuẫn trong cách trình bày.

\textbf{Hạn chế của báo cáo.} Báo cáo không có dữ liệu gốc của bài báo, nên phần tái hiện chỉ kiểm tra được tính nhất quán của số liệu đã công bố và hành vi của phương pháp trên dữ liệu mô phỏng, chứ không khẳng định được kết quả trên dữ liệu thật. Các thông số của dự án hệ thống thông tin mẫu là giả định để minh họa.

\textbf{Hướng phát triển.} Nếu nhận được dữ liệu gốc từ tác giả, có thể chạy lại mô hình để kiểm tra trực tiếp các nghi vấn đã nêu. Ngoài ra, có thể dùng bộ câu hỏi của bài báo để khảo sát các dự án tại Việt Nam, hoặc điều chỉnh cho dự án hệ thống thông tin, rồi so sánh với kết quả ở Yemen. Đó cũng là hướng mà tác giả gợi ý ở mục 6.2 của bài báo.
